% !TeX root = ../navier-stokes-manuscript.tex
% ============================================================
% 2.5 Simultaneous Sobolev Intersections
% ============================================================

\subsection{Simultaneous Sobolev Intersections}\label{sub:SimultaneousSobolevIntersections}

The word smooth can now be connected to the tower one derivative order at a time.
In three spatial dimensions, Sobolev embedding gives
\[
  H^m(\R^3)
  \hookrightarrow
  C^k(\R^3)
  \qquad
  \text{whenever }m>k+\frac32.
\]
Fix any requested spatial derivative order \(k\).
One finite choice of \(m>k+3/2\) makes every derivative of \(u\) through order \(k\) continuous and bounded.
Allowing \(m\) to range through all finite orders gives
\[
  \bigcap_{m=0}^{\infty}H^m(\R^3)
  \hookrightarrow
  C^\infty(\R^3).
\]
No individual Sobolev row has been called infinitely smooth.
The intersection supplies every finite row, and each desired derivative chooses the one finite row high enough to control it.

The temporal direction has the same form on a bounded strict interval \(I=(0,T)\) with \(T<T_*\).
For a Hilbert space \(X\), the one-dimensional Sobolev embedding gives
\[
  H^r(I;X)
  \hookrightarrow
  C^q(\overline I;X)
  \qquad
  \text{whenever }r>q+\frac12.
\]
Fix a requested temporal derivative order \(q\), then choose one finite \(r>q+1/2\).
Allowing all finite temporal orders gives
\[
  \bigcap_{r=0}^{\infty}H^r(I;X)
  \hookrightarrow
  C^\infty(\overline I;X).
\]
The temporal intersection has its own intersection-to-smoothness argument, just as the spatial intersection does.

The mixed jet asks for both choices at once.
Consider the coordinatewise mixed intersection
\[
  \mathscr J_\infty(I)
  :=
  \bigcap_{r=0}^{\infty}
  \bigcap_{m=0}^{\infty}
  H^r\bigl(I;H^m(\R^3)\bigr).
\]
Choose any temporal order \(q\) and spatial order \(k\).
Selecting finite indices \(r>q+1/2\) and \(m>k+3/2\) gives
\[
  H^r\bigl(I;H^m(\R^3)\bigr)
  \hookrightarrow
  C^q\bigl(\overline I;C^k(\R^3)\bigr).
\]
Since \(q\) and \(k\) were arbitrary, the full mixed intersection satisfies
\[
  \mathscr J_\infty(I)
  \hookrightarrow
  C^\infty\bigl(\overline I\times\R^3\bigr).
\]
Every mixed derivative \(\partial_t^q\partial_x^\alpha u\) is reached by a finite coordinate of the intersection.

These intersections are simultaneous readings of the fluid.
The spatial column reads all spatial heights of \(u(t)\).
The temporal column reads all time responses of a fixed field-valued coordinate.
The critical-height column reads the scalar \(H\) and the nested energies and production terms exposed by its triangular identity.
The mixed intersection reads every \(J_{m,\alpha}\) together with the pressure coordinates and recurrence that relate them.
The columns do not describe successive configurations through which the fluid must pass.
They expose different aspects of the same coupled equation at the same time.

The critical-height column also has a strict limitation.
The scalar \(H(t)\) records the size of \(\Lambda^{1/2}u(t)\), so it forgets the field's orientation, phase, and spatial arrangement.
Even complete temporal information about \(H\) cannot reconstruct \(u\) by itself.
Its value comes from the exact identities that place its derivatives on the compatible graph of the full mixed jet.
The velocity and pressure coordinates remain the object that identifies which spatial field produced the scalar readout.

Intersection means coordinatewise membership.
For the mixed object, the statement
\[
  u\in\mathscr J_\infty(I)
\]
means that for every finite pair \((r,m)\),
\[
  \norm{u}_{H^r(I;H^m)}<\infty.
\]
It places no demand on the infinite series
\[
  \sum_{r,m=0}^{\infty}
  \norm{u}_{H^r(I;H^m)}^2.
\]
That series may diverge while every finite coordinate needed for smoothness remains finite.
Different temporal and spatial coordinates may also carry different constants.
The smoothness implication uses one finite coordinate after the desired derivative order has been chosen.

This distinction changes the proof design.
There is no need to dominate the entire infinite tower by one universal all-orders norm.
The work can be done at arbitrary finite temporal and spatial depth, provided the finite views agree whenever they contain the same derivative.
The smallest useful view must preserve two independent cutoffs and the relation between a temporal row and its derivative.
That requirement produces the adjacent-row rectangle.
